% Density revision 2026-10-04: answers below questions; compact related content.
% 4.x ছবি in HTML
\section{HTML-এ ছবি যুক্ত করা}

% Composition pass 2026-10-04
কলেজের পরিচিতি বা দোকানের পণ্য শুধু লেখা দিয়ে দেখালে চেনা কঠিন হতে পারে।

\chFourMotivation{দৃশ্যমান পরিচয় $\rightarrow$ img $\rightarrow$ src থেকে ছবির রিসোর্স লোড ও alt দিয়ে বিকল্প অর্থ $\rightarrow$ ক্যাম্পাসের ছবি, ব্যানার ও পণ্য প্রদর্শন।}


\begin{learningbox}
এই অংশ শেষে শিক্ষার্থী পারবে:
\begin{itemize}
\item HTML পেজে ছবি যুক্ত করার নিয়ম ব্যাখ্যা করতে।
\item \texttt{<img>} ট্যাগ এম্পটি ট্যাগ কেন তা লিখতে।
\item \texttt{src}, \texttt{alt}, \texttt{width} ও \texttt{height} অ্যাট্রিবিউট ব্যবহার করতে।
\item ছবির অবস্থানের পথ, ছবির ফরম্যাট, ছবির অবস্থান, বর্ডার ও ছবিতে লিংক বুঝতে।
\item ধারকের মাপ অনুযায়ী ছবি ইঙ্গিত যেমন \texttt{srcset}, \texttt{sizes} ও \texttt{loading} চিনতে।
\item সারণির ঘরের ভিতরে ছবি বসানোর কোড লিখতে।
\item HSC ব্যবহারিক অনুশীলন ও বোর্ড পরীক্ষার ছবি সম্পর্কিত প্রশ্নের উত্তর দিতে।
\end{itemize}
\end{learningbox}

হাইপারলিংক শেখার পর ছবি শেখা সহজ, কারণ ওয়েবপেজে ছবি শুধু দেখানোই নয়, ছবির ওপর লিংকও দেওয়া যায়। কোনো কলেজের ওয়েবসাইটে লোগো, ব্যানার, শিক্ষকের ছবি, ক্যাম্পাসের ছবি, বিজ্ঞপ্তি বা ছবির সংগ্রহ দেখাতে HTML-এর \texttt{<img>} ট্যাগ ব্যবহার করা হয়।

\subsection{ছবি ট্যাগ কী?}

\begin{definitionbox}{ছবির ট্যাগ (Image Tag)}
ওয়েবপেজে ছবি যুক্ত করতে \texttt{<img>} ট্যাগ ব্যবহার করা হয়। এটি ছবির ট্যাগ।
\end{definitionbox}

তবে শুধু \texttt{<img>} লিখলে ছবি দেখা যাবে না। ব্রাউজারকে ছবির ফাইল কোথায় আছে তা জানাতে \texttt{src} অ্যাট্রিবিউট ব্যবহার করতে হয়। ছবির বিকল্প লেখা দেওয়ার জন্য \texttt{alt} অ্যাট্রিবিউট ব্যবহার করা হয়। ছবি দেখানো না গেলে এই লেখা দেখা যেতে পারে; স্ক্রিন রিডারও এটি ব্যবহার করে।

\subsection{\texttt{<img>} কেন এম্পটি ট্যাগ?}

\texttt{<img>} ওয়েবপেজে ছবি যুক্ত করে। এটি ভয়েড এলিমেন্টের শুরুর ট্যাগ; পাঠ্যবইয়ে এম্পটি ট্যাগ নামে পরিচিত।

\texttt{<img>} ট্যাগকে এম্পটি ট্যাগ বলা হয়, কারণ:

\begin{itemize}
\item এর শেষের ট্যাগ নেই।
\item এর ভিতরে লেখা বিষয়বস্তু থাকে না।
\item এটি অ্যাট্রিবিউট ব্যবহার করে ছবি প্রদর্শন করে।
\end{itemize}



\subsection{ছবি যুক্ত করার সাধারণ রূপ}

\begin{examplebox}{কোড}
\par\noindent\begin{minipage}{\linewidth}
\begin{verbatim}
<img src="image-name">
\end{verbatim}
\end{minipage}\par
\end{examplebox}

আরও মানসম্মত সিনট্যাক্স:

\begin{examplebox}{কোড}
\par\noindent\begin{minipage}{\linewidth}
\begin{verbatim}
<img src="image-name.jpg" alt="Image description">
\end{verbatim}
\end{minipage}\par
\end{examplebox}

\subsubsection{উদাহরণ}

\begin{examplebox}{কোড}
\par\noindent\begin{minipage}{\linewidth}
\begin{verbatim}
<img src="flower.jpg" alt="Flower">
\end{verbatim}
\end{minipage}\par
\end{examplebox}

\begin{center}
\begin{tabular}{|p{0.25\textwidth}|p{0.55\textwidth}|}
\hline
\textbf{অংশ} & \textbf{কাজ} \\
\hline
\texttt{<img>} & ছবি ট্যাগ \\
\hline
\texttt{src} & ছবির ফাইলের পথ বা নাম নির্দেশ করে \\
\hline
\texttt{"flower.jpg"} & ছবির ফাইলের নাম \\
\hline
\texttt{alt} & ছবি না দেখালে বিকল্প লেখা দেয় \\
\hline
\texttt{"Flower"} & বিকল্প লেখা \\
\hline
\end{tabular}
\end{center}

\begin{examplebox}{ফল ও ব্যাখ্যা}
\textbf{ফল:} ব্রাউজারে \texttt{flower.jpg} ছবি দেখা যাবে। ছবি লোড না হলে \textbf{Flower} লেখা দেখা যেতে পারে।

\textbf{ব্যাখ্যা:} \texttt{src} ছবির ফাইলের অবস্থান জানায় এবং \texttt{alt} ছবির বিকল্প লেখা দেয়।
\end{examplebox}

\begin{center}
\begin{tikzpicture}[
  part/.style={draw, rounded corners, fill=green!6, minimum height=0.72cm, align=center, font=\scriptsize},
  label/.style={draw, rounded corners, fill=blue!6, align=center, font=\scriptsize},
  arrow/.style={->, thick}
]
\node[part, minimum width=1.25cm] (tag) {\texttt{<img}};
\node[part, minimum width=2.65cm, right=0cm of tag] (src) {\texttt{src="flower.jpg"}};
\node[part, minimum width=2.55cm, right=0cm of src] (alt) {\texttt{alt="Flower"}};
\node[part, minimum width=0.7cm, right=0cm of alt] (end) {\texttt{>}};
\node[label, below=0.7cm of tag] (ltag) {image tag\\empty tag};
\node[label, below=0.7cm of src] (lsrc) {image file-এর\\path/name};
\node[label, below=0.7cm of alt] (lalt) {image না দেখালে\\alternative text};
\draw[arrow] (ltag) -- (tag);
\draw[arrow] (lsrc) -- (src);
\draw[arrow] (lalt) -- (alt);
\end{tikzpicture}\\[3pt]
\textit{চিত্র: \texttt{<img>} ট্যাগের গুরুত্বপূর্ণ অংশ}
\end{center}

\subsection{\texttt{src} অ্যাট্রিবিউট}

\begin{definitionbox}{\texttt{src}}
\texttt{src} অ্যাট্রিবিউট ছবির ফাইলের ঠিকানা জানায়। ব্রাউজার ওই ঠিকানা থেকে ছবিটি সংগ্রহ করে।
\end{definitionbox}

\subsubsection{উদাহরণ}

\begin{chFourWholeExample}{কোড ও ব্রাউজারে ফল}
\noindent
\begin{minipage}[t]{.50\linewidth}\vspace{0pt}

\begin{verbatim}
<img src="rose.jpg" alt="Rose">
\end{verbatim}

\end{minipage}\hfill
\begin{minipage}[t]{.46\linewidth}\vspace{0pt}
\textbf{ফল:} ব্রাউজার একই ফোল্ডার থেকে \texttt{rose.jpg} ছবি খুঁজে দেখাবে।
\end{minipage}
\chFourWhy{HTML ফাইল ও ছবি একই ফোল্ডারে থাকলে শুধু ছবির ফাইলের নাম লিখলেই চলে। অন্য ফোল্ডারে থাকলে পথ দিতে হয়।}
\end{chFourWholeExample}



\subsection{\texttt{alt} অ্যাট্রিবিউট}

\begin{definitionbox}{\texttt{alt}}
\texttt{alt} অ্যাট্রিবিউট ছবির বিকল্প লেখা নির্ধারণ করে। ছবি লোড না হলে বা স্ক্রিন রিডার ব্যবহার করলে এই লেখা কাজে লাগে।
\end{definitionbox}

\subsubsection{উদাহরণ}

% Split revision: code and its explanation share one row.
\par\noindent
\begin{minipage}[t]{.54\linewidth}\vspace{0pt}\small
\textbf{\color{chFourTeal}কোড}\par
\begin{verbatim}
<img src="college.jpg" alt="ABC College Building">
\end{verbatim}

\par\smallskip
\end{minipage}\hfill
\begin{minipage}[t]{.42\linewidth}\vspace{0pt}\small
\textbf{\color{chFourTeal}ফল ও ব্যাখ্যা}\par
\textbf{ফল:} \texttt{college.jpg} ছবি লোড হলে কলেজ ভবনের ছবি দেখা যাবে। ছবি লোড না হলে \textbf{ABC College Building} লেখা দেখা যেতে পারে।

\textbf{ব্যাখ্যা:} ছবির উদ্দেশ্য অনুযায়ী বিকল্প লেখা দিলে ছবি দেখতে না পাওয়া ব্যবহারকারীও তথ্য বুঝতে পারেন। কেবল সাজানোর জন্য ব্যবহৃত ছবিতে \texttt{alt=""} দিয়ে তা বোঝানো যায়।
\par\smallskip
\end{minipage}\par





\subsection{\texttt{width} ও \texttt{height} অ্যাট্রিবিউট}

\begin{definitionbox}{\texttt{width} ও \texttt{height}}
\texttt{width} ছবির প্রস্থ এবং \texttt{height} ছবির উচ্চতা নির্ধারণ করে।
\end{definitionbox}

\subsubsection{সিনট্যাক্স}

\begin{examplebox}{কোড}
\par\noindent\begin{minipage}{\linewidth}
\begin{verbatim}
<img src="image.jpg" width="300" height="200" alt="Image">
\end{verbatim}
\end{minipage}\par
\end{examplebox}

\subsubsection{উদাহরণ}

% Split revision: code and its explanation share one row.
\par\noindent
\begin{minipage}[t]{.54\linewidth}\vspace{0pt}\small
\textbf{\color{chFourTeal}কোড}\par
\begin{verbatim}
<img src="flower.png" width="400" height="200" alt="Flower">
\end{verbatim}

\par\smallskip
\end{minipage}\hfill
\begin{minipage}[t]{.42\linewidth}\vspace{0pt}\small
\textbf{\color{chFourTeal}ফল ও ব্যাখ্যা}\par
\textbf{ফল:} ব্রাউজার \texttt{flower.png} ছবি 400 পিক্সেল প্রস্থ এবং 200 পিক্সেল উচ্চতায় দেখাবে।

\textbf{ব্যাখ্যা:} \texttt{width} ও \texttt{height} অ্যাট্রিবিউট ছবির প্রদর্শিত আকার নিয়ন্ত্রণ করে। শুধু প্রস্থ দিলে ব্রাউজার স্বাভাবিকভাবে অনুপাত ধরে উচ্চতা নির্ধারণ করতে পারে।
\par\smallskip
\end{minipage}\par


\begin{keypointbox}{আধুনিক HTML-এর নিয়ম}
ধারকের মাপ অনুযায়ী বদলানো নকশার জন্য CSS ব্যবহার করা ভালো। \texttt{width} ও \texttt{height} দিয়ে ছবির প্রদর্শিত প্রস্থ ও উচ্চতাও নির্ধারণ করা যায়।
\end{keypointbox}

\subsection{ছবির ফাইলের ফরম্যাট}

ওয়েবপেজে \texttt{.jpg}, \texttt{.jpeg}, \texttt{.png}, \texttt{.gif} ও \texttt{.bmp} ফরম্যাটের ছবি ব্যবহার করা যায়। \texttt{.svg}, \texttt{.webp} ও \texttt{.avif}-ও আধুনিক ওয়েবে ব্যবহৃত হয়। ফরম্যাটভেদে ফাইলের আকার, স্বচ্ছতা ও সংকোচনের ধরন আলাদা হয়।

\begin{center}
\begin{tabular}{|p{0.25\textwidth}|p{0.55\textwidth}|}
\hline
\textbf{ফরম্যাট} & \textbf{ব্যবহার} \\
\hline
\texttt{.jpg} / \texttt{.jpeg} & ফটোগ্রাফ ও ক্যামেরায় তোলা ছবি \\
\hline
\texttt{.png} & স্বচ্ছ পটভূমির ছবি, স্ক্রিনশট ও লোগো \\
\hline
\texttt{.gif} & সহজ অ্যানিমেশন বা ছোট গ্রাফিক \\
\hline
\texttt{.svg} & ভেক্টর ছবি, লোগো ও আইকন \\
\hline
\texttt{.bmp} & বিটম্যাপ ছবি; ওয়েবে কম ব্যবহৃত \\
\hline
\texttt{.webp} & সংকুচিত স্থির বা অ্যানিমেটেড ছবি \\
\hline
\texttt{.avif} & বেশি সংকোচনের সুবিধাযুক্ত ছবি \\
\hline
\end{tabular}
\end{center}

HTML মানে একটি নির্দিষ্ট ছবির ফরম্যাট বাধ্যতামূলক নয়; ব্রাউজারের সমর্থন ভিন্ন হতে পারে। WebP ও AVIF-এ সংকুচিত ছবি রাখা যায়, আর SVG ভেক্টর ছবির আকার পরিবর্তনেও রেখা স্পষ্ট রাখতে পারে। ছবির ধরন, আকার ও প্রয়োজনীয় ব্রাউজারের সমর্থন দেখে ফরম্যাট বেছে নিতে হয়।

\subsection{ছবির অবস্থানের পথ}

ছবির ফাইল কোথায় রাখা আছে তার উপর \texttt{src} মান নির্ভর করে।

\subsubsection{একই ফোল্ডারে ছবি থাকলে}

\begin{chFourWholeExample}{কোড ও ব্রাউজারে ফল}
\noindent
\begin{minipage}[t]{.50\linewidth}\vspace{0pt}

\begin{verbatim}
<img src="flower.jpg" alt="Flower">
\end{verbatim}

\end{minipage}\hfill
\begin{minipage}[t]{.46\linewidth}\vspace{0pt}
HTML ফাইল ও \texttt{flower.jpg} একই ফোল্ডারে থাকলে ব্রাউজার ছবি দেখাবে।
\end{minipage}

\end{chFourWholeExample}

\subsubsection{ছবি আলাদা ফোল্ডারে থাকলে}

\begin{chFourWholeExample}{কোড ও ব্রাউজারে ফল}
\noindent
\begin{minipage}[t]{.50\linewidth}\vspace{0pt}

\begin{verbatim}
<img src="images/flower.jpg" alt="Flower">
\end{verbatim}

\end{minipage}\hfill
\begin{minipage}[t]{.46\linewidth}\vspace{0pt}
এখানে ব্রাউজার \texttt{images} ফোল্ডারের ভিতরে \texttt{flower.jpg} ছবি খুঁজবে।
\end{minipage}

\end{chFourWholeExample}

\subsubsection{অনলাইন ছবি হলে}

% Split revision: code and its explanation share one row.
\par\noindent
\begin{minipage}[t]{.54\linewidth}\vspace{0pt}\small
\textbf{\color{chFourTeal}কোড}\par
\begin{verbatim}
<img src="https://example.com/flower.jpg" alt="Flower">
\end{verbatim}

\par\smallskip
\end{minipage}\hfill
\begin{minipage}[t]{.42\linewidth}\vspace{0pt}\small
\textbf{\color{chFourTeal}ফল ও ব্যাখ্যা}\par
ব্রাউজার ইন্টারনেট থেকে ছবি লোড করার চেষ্টা করবে। অনলাইন ছবি ব্যবহার করলে ঠিকানা ঠিক থাকতে হবে এবং ইন্টারনেট সংযোগ দরকার হতে পারে।
\par\smallskip
\end{minipage}\par


\begin{mistakebox}{ছবি না দেখালে কী যাচাই করবে?}
\begin{itemize}
\item \texttt{src} পথ ভুল কি না।
\item ছবির ফাইলের নাম ঠিক আছে কি না।
\item এক্সটেনশন \texttt{.jpg}, \texttt{.png}, \texttt{.gif} ঠিক আছে কি না।
\item ছবির ফাইল নির্ধারিত ফোল্ডারে আছে কি না।
\item ইংরেজি বড় বা ছোট হাতের অক্ষর ও বানানে অমিল আছে কি না।
\end{itemize}
\end{mistakebox}

\Needspace{7\baselineskip}
\subsection{ছবির আকার নিয়ন্ত্রণ}

\subsubsection{উদাহরণ: 300 পিক্সেল প্রস্থ ও 200 পিক্সেল উচ্চতার ছবি}

% Split revision: code and its explanation share one row.
\par\noindent
\begin{minipage}[t]{.54\linewidth}\vspace{0pt}\small
\textbf{\color{chFourTeal}কোড}\par
\begin{verbatim}
<img src="nature.jpg" width="300" height="200" alt="Nature">
\end{verbatim}

\par\smallskip
\end{minipage}\hfill
\begin{minipage}[t]{.42\linewidth}\vspace{0pt}\small
\textbf{\color{chFourTeal}ফল ও ব্যাখ্যা}\par
ব্রাউজার \texttt{nature.jpg} ছবি 300 পিক্সেল প্রস্থ এবং 200 পিক্সেল উচ্চতায় দেখাবে।
\par\smallskip
\end{minipage}\par


\subsubsection{পর্দার মাপ অনুযায়ী CSS-এর রূপ}

% Split revision: code and its explanation share one row.
\par\noindent
\begin{minipage}[t]{.54\linewidth}\vspace{0pt}\small
\textbf{\color{chFourTeal}কোড}\par
\begin{verbatim}
<img src="nature.jpg" alt="Nature"
     style="width:100%; max-width:500px; height:auto;">
\end{verbatim}

\par\smallskip
\end{minipage}\hfill
\begin{minipage}[t]{.42\linewidth}\vspace{0pt}\small
\textbf{\color{chFourTeal}ফল ও ব্যাখ্যা}\par
ছবিটি তার ধারক অংশের প্রস্থ অনুযায়ী ছোট-বড় হবে, তবে 500 CSS পিক্সেলের বেশি চওড়া হবে না। ধারক অংশের মাপ সাধারণত পর্দার মাপের সঙ্গে বদলায়। \texttt{height:auto} ছবির প্রস্থ-উচ্চতার অনুপাত ধরে রাখে।
\par\smallskip
\end{minipage}\par


\subsection{পর্দা অনুযায়ী ছবি বাছাইয়ের আধুনিক ইঙ্গিত}

% Split revision: code and its explanation share one row.
\par\noindent
\begin{minipage}[t]{.54\linewidth}\vspace{0pt}\small
\textbf{\color{chFourTeal}কোড}\par
\begin{verbatim}
<img src="college-400.jpg"
     srcset="college-400.jpg 400w, college-800.jpg 800w"
     sizes="(max-width: 600px) 100vw, 50vw"
     loading="lazy"
     decoding="async"
     width="400" height="250"
     alt="ABC College Building">
\end{verbatim}

\par\smallskip
\end{minipage}\hfill
\begin{minipage}[t]{.42\linewidth}\vspace{0pt}\small
\textbf{\color{chFourTeal}ফল ও ব্যাখ্যা}\par
ব্রাউজার দেখার এলাকার মাপ, ছবির প্রদর্শিত মাপ ও যন্ত্রের পিক্সেলের ঘনত্ব বিবেচনা করে উপযুক্ত ফাইল বেছে নিতে পারে। \texttt{loading="lazy"} দিয়ে প্রয়োজনের কাছাকাছি সময়ে ছবি লোড করার ইঙ্গিত দেওয়া হয়; \texttt{decoding="async"} ছবি ডিকোডের জন্য অন্য বিষয়বস্তুর প্রদর্শন অপেক্ষায় না রাখার ইঙ্গিত দেয়। \texttt{width} ও \texttt{height} আগে থেকে জায়গা রাখতে সাহায্য করে।
\par\smallskip
\end{minipage}\par


\subsection{ছবির অবস্থান}

পরীক্ষায় পরিচিত পুরোনো HTML-এ ছবির অবস্থান নির্ধারণের জন্য \texttt{align} অ্যাট্রিবিউট দেখা যায়। তবে আধুনিক ওয়েব নকশায় CSS ব্যবহার করা বেশি মানসম্মত।


\textbf{ছবিটি আসলে কোথায় থাকে?} এই পাঠের সাধারণ উদাহরণে \textbf{src} ছবির ফাইলের ঠিকানা দেয়; HTML ফাইলে ছবির সব বাইট লিখে রাখা হয় না। ব্রাউজার পথটি বর্তমান পেজের অবস্থানের সঙ্গে মিলিয়ে ফাইল খুঁজে লোড করে। তাই HTML পাঠানোর পাশাপাশি সংশ্লিষ্ট ছবিও সঠিক ফোল্ডারে রাখতে হবে। \textbf{alt} ছবির কাজ বা প্রয়োজনীয় তথ্য বোঝাবে; নিছক সাজসজ্জার ছবিতে \texttt{alt=""} উপযুক্ত হতে পারে।


\subsubsection{পরীক্ষায় পরিচিত উদাহরণ}

\begin{examplebox}{কোড}
\par\noindent\begin{minipage}{\linewidth}
\begin{verbatim}
<img src="z1.jpg" width="200" height="150"
     align="left" alt="Sample Image">
\end{verbatim}
\end{minipage}\par
\end{examplebox}

\begin{center}
\begin{tabular}{|p{0.25\textwidth}|p{0.55\textwidth}|}
\hline
\textbf{align-এর মান} & \textbf{কাজ} \\
\hline
\texttt{left} & ছবি বাম পাশে রাখে \\
\hline
\texttt{right} & ছবি ডান পাশে রাখে \\
\hline
\texttt{top} & আশপাশের লেখার ওপরের অংশের সঙ্গে অবস্থান নির্দেশ করে \\
\hline
\texttt{middle} & আশপাশের লেখার মাঝের অংশের সঙ্গে অবস্থান নির্দেশ করে \\
\hline
\texttt{bottom} & আশপাশের লেখার নিচের অংশের সঙ্গে অবস্থান নির্দেশ করে \\
\hline
\end{tabular}
\end{center}

\subsubsection{আধুনিক CSS-এর রূপ}

\begin{examplebox}{কোড}
\par\noindent\begin{minipage}{\linewidth}
\begin{verbatim}
<img src="z1.jpg" alt="Sample Image"
     style="float:left; width:200px; height:150px;">
\end{verbatim}
\end{minipage}\par
\end{examplebox}

\subsection{ছবির বর্ডার}

পাঠ্যবইয়ের পুরোনো রীতির কোডে ছবির চারপাশে বর্ডার দিতে \texttt{border} অ্যাট্রিবিউট দেখা যায়।

\subsubsection{পরীক্ষায় পরিচিত উদাহরণ}

\begin{chFourWholeExample}{কোড ও ব্রাউজারে ফল}
\noindent
\begin{minipage}[t]{.50\linewidth}\vspace{0pt}

\begin{verbatim}
<img src="z1.jpg" width="200" height="150"
     border="1" alt="Sample Image">
\end{verbatim}

\end{minipage}\hfill
\begin{minipage}[t]{.46\linewidth}\vspace{0pt}
ছবির চারপাশে 1 পিক্সেল বর্ডার দেখা যাবে।
\end{minipage}
\chFourWhy{\textbf{src} থেকে ছবি লোড হয়; \textbf{alt} সেই ছবির বিকল্প অর্থ দেয়। ফাইলের পথ সঠিক হতে হবে।}
\end{chFourWholeExample}

\subsubsection{আধুনিক CSS-এর রূপ}

\begin{examplebox}{কোড}
\par\noindent\begin{minipage}{\linewidth}
\begin{verbatim}
<img src="z1.jpg" alt="Sample Image"
     style="border:1px solid black;">
\end{verbatim}
\end{minipage}\par
\end{examplebox}

\subsection{ছবিতে লিংক}

ছবিতে ক্লিক করে অন্য পেজ বা ওয়েবসাইট খোলার ব্যবস্থা করতে চাইলে \texttt{<a>} ট্যাগের ভিতরে \texttt{<img>} ট্যাগ বসাতে হয়।

\subsubsection{সিনট্যাক্স}

\begin{examplebox}{কোড}
\par\noindent\begin{minipage}{\linewidth}
\begin{verbatim}
<a href="url">
  <img src="image.jpg" alt="Image">
</a>
\end{verbatim}
\end{minipage}\par
\end{examplebox}

\subsubsection{উদাহরণ}

% Split revision: code and its explanation share one row.
\par\noindent
\begin{minipage}[t]{.54\linewidth}\vspace{0pt}\small
\textbf{\color{chFourTeal}কোড}\par
\begin{verbatim}
<a href="https://www.example.com">
  <img src="banner.jpg" width="200" height="95"
     alt="Visit Example website">
</a>
\end{verbatim}

\par\smallskip
\end{minipage}\hfill
\begin{minipage}[t]{.42\linewidth}\vspace{0pt}\small
\textbf{\color{chFourTeal}ফল ও ব্যাখ্যা}\par
ব্রাউজারে ব্যানারের ছবি দেখা যাবে। ছবিতে ক্লিক করলে নির্দিষ্ট ওয়েবসাইট খুলবে।
\par\smallskip
\end{minipage}\par


\begin{keypointbox}{ছবিতে লিংকের নিয়ম}
ছবিতে লিংকের \texttt{alt} লেখা ছবির চেহারার বর্ণনার চেয়ে লিংকের উদ্দেশ্য বা গন্তব্য বোঝানো ভালো।
\end{keypointbox}



\subsubsection{ছবির ফরম্যাট বেছে নেওয়ার কারণ}
ক্যাম্পাসের ফটোগ্রাফে JPEG-এ ফাইল ছোট করা যায়, তবে অতিরিক্ত সংকোচনে সূক্ষ্ম তথ্য হারাতে পারে। স্বচ্ছ পটভূমির লোগো বা ধারালো লেখা-যুক্ত স্ক্রিনশটে PNG উপযোগী; এতে তথ্য অক্ষুণ্ণ রাখা সংকোচন (lossless compression) হয়। সাধারণ অ্যানিমেশনে GIF ব্যবহার করা যায়, কিন্তু প্রতি ফ্রেমে সীমিত রঙের কারণে সব ফটোগ্রাফের জন্য এটি ভালো পছন্দ নয়। রেখাভিত্তিক লোগোতে SVG বড় করলেও রেখা স্পষ্ট রাখতে পারে। অর্থাৎ শুধু ফাইলের এক্সটেনশন নয়—ছবির তথ্য, স্বচ্ছতা, মান ও ফাইলের মাপ দেখে সিদ্ধান্ত নিতে হবে।

\subsubsection{থাম্বনেইল থেকে বড় ছবিতে যাওয়া}
\textbf{থাম্বনেইল (Thumbnail)} হলো বড় ছবির ছোট পূর্বরূপ। সংগ্রহের পেজে ছোট ছবি দেখে প্রয়োজনীয় বড় ছবিটি খোলা সুবিধাজনক। নিচে দুই ফাইলই নিজে তৈরি করে একই ফোল্ডারে রাখতে হবে; ছোটটির অনুপাত 4:3 ধরা হয়েছে।
\par\noindent\begin{minipage}[t]{.50\linewidth}\vspace{0pt}\small
\textbf{\color{chFourTeal}ছবি ও লিংকের সমন্বয়}\par
\begin{verbatim}
<a href="library-large.jpg">
  <img src="library-small.jpg"
       alt="College library"
       width="128" height="96">
</a>
\end{verbatim}
\end{minipage}\hfill
\begin{minipage}[t]{.46\linewidth}\vspace{0pt}\small
\textbf{\color{chFourTeal}ফল ও ব্যাখ্যা}\par
ছোট ছবিটি 128\,$\times$\,96 পিক্সেলে দেখা যাবে। সেটি নির্বাচন করলে বড় ফাইলটি খুলবে। \texttt{src} দেখানোর ছোট ছবি, \texttt{href} যাওয়ার বড় ছবি। আলাদা ছোট ফাইল রাখলে প্রথম পেজে বড় ফাইলটি নামানোর প্রয়োজন কমে।
\end{minipage}\par

\subsubsection{তথ্যের ছবি ও পটভূমির ছবি}
ব্যানারে ভর্তি-তারিখ বা কলেজের নাম থাকলে তা পাঠের তথ্য; তখন \texttt{img} ও উপযুক্ত \texttt{alt} দেবে। শুধু সাজসজ্জার পটভূমি পটভূমিতে দিতে CSS-এর \texttt{background-image} ব্যবহার করা যায়। পুরোনো \texttt{body background} রীতির আধুনিক বিকল্প এটি।
\par\noindent\begin{minipage}[t]{.50\linewidth}\vspace{0pt}\small
\textbf{\color{chFourTeal}body-এর CSS রীতি}\par
\begin{verbatim}
body {
  background-color: white;
  background-image:
    url("paper-texture.png");
}
\end{verbatim}
\end{minipage}\hfill
\begin{minipage}[t]{.46\linewidth}\vspace{0pt}\small
\textbf{\color{chFourTeal}কোথায় লিখবে, কী দেখা যাবে}\par
কোডটি \texttt{head}-এর \texttt{style} এলিমেন্টে লিখবে। ফাইল থাকলে সাদা পটভূমির ওপর ছবিটি স্বাভাবিক রীতিতে পুনরাবৃত্ত হবে; না এলে সাদা রং থাকবে। CSS-এর background-এ \texttt{alt} নেই; প্রয়োজনীয় তথ্য সেখানে একা রাখবে না।
\end{minipage}\par
ভর্তি-শেষের তারিখ শুধু পটভূমির ছবিতে থাকলে ছবি না আসা বা সহায়ক পাঠযন্ত্র ব্যবহারের সময় তথ্যটি হারাতে পারে। তাই প্রয়োজনীয় তারিখ পাঠ হিসেবেও রাখা উচিত।


\subsubsection{ছবির ফরম্যাটের পূর্ণরূপ}
JPEG-এর পূর্ণরূপ Joint Photographic Experts Group; প্রচলিত ফাইল এক্সটেনশন \texttt{.jpg} ও \texttt{.jpeg}। PNG হলো Portable Network Graphics এবং GIF হলো Graphics Interchange Format; এক্সটেনশন যথাক্রমে \texttt{.png} ও \texttt{.gif}। ফাইলের নামের এক্সটেনশন বদলালেই ছবির প্রকৃত ফরম্যাট বদলে যায় না; ফরম্যাট বদলাতে উপযুক্তভাবে রূপান্তর ও সংরক্ষণ করতে হয়।

\subsubsection{পুরোনো ছবির অ্যাট্রিবিউটে ফাঁক}
পুরোনো HTML-এ \texttt{hspace} ছবির বাম ও ডান পাশে এবং \texttt{vspace} ওপর ও নিচে ফাঁকের পরিমাণ নির্দেশ করত। যেমন, \texttt{hspace="10"} দুই পাশে 10 পিক্সেল করে ফাঁক বোঝায়। এগুলো ছবির \texttt{width} বা \texttt{height} বদলানোর অ্যাট্রিবিউট নয়। নতুন কোডে ছবির বাইরের ফাঁক CSS-এর \texttt{margin} দিয়ে নির্ধারণ করা উচিত।


\subsection{সারণির ভিতরে ছবি}

HTML সারণির ঘরের ভিতরেও ছবি রাখা যায়। ফলের সারণি, পণ্যের সারণি, ব্যক্তির পরিচিতির সারণি বা ছবির সংগ্রহের বিন্যাসে এটি কাজে লাগে।

\subsubsection{উদাহরণ}

% Split revision: code and its explanation share one row.
\par\noindent
\begin{minipage}[t]{.54\linewidth}\vspace{0pt}\small
\textbf{\color{chFourTeal}কোড}\par
\begin{verbatim}
<table border="1">
  <tr>
    <td><img src="rose.gif" alt="Rose" width="100"></td>
    <td>Rose Image</td>
  </tr>
</table>
\end{verbatim}

\par\smallskip
\end{minipage}\hfill
\begin{minipage}[t]{.42\linewidth}\vspace{0pt}\small
\textbf{\color{chFourTeal}ফল ও ব্যাখ্যা}\par
সারণির প্রথম ঘরে গোলাপের ছবি এবং দ্বিতীয় ঘরে \textbf{Rose Image} লেখা দেখা যাবে।
\par\smallskip
\end{minipage}\par


\subsection{ছবি যুক্ত করার পূর্ণাঙ্গ উদাহরণ}

\Needspace{6\baselineskip}
\par\noindent{\color{chFourTeal}\bfseries\small কোড: বাম কলাম শেষে ডান কলাম পড়বে}\par
\begin{multicols}{2}
\begin{verbatim}
<!DOCTYPE html>
<html>
<head>
  <meta charset="UTF-8">
  <title>Image Example</title>
</head>
<body>
  <h1>My College</h1>
  <p>Below is an image of our college:</p>
  <img src="college.jpg"
       alt="ABC College Building"
       width="400" height="250">
  <p>Click the banner below to visit our website:</p>
  <a href="https://www.example.com">
    <img src="banner.jpg"
         alt="College Website Banner"
         width="300" height="100">
  </a>
</body>
</html>
\end{verbatim}
\end{multicols}


\textbf{\color{chFourTeal}ফল কী হবে?}\par
ব্রাউজারে দেখা যাবে:
\begin{itemize}
\item \textbf{My College} শিরোনাম।
\item একটি কলেজ ছবি।
\item একটি ব্যানারের ছবি।
\item ব্যানারের ছবিতে ক্লিক করলে নির্দিষ্ট ওয়েবসাইট খুলবে।
\end{itemize}
\par\smallskip



\subsection{সাধারণ ভুল}

\begin{multicols}{2}
\small
\begin{itemize}
\item \texttt{<image>} ট্যাগ লেখা; সঠিক ট্যাগ হলো \texttt{<img>}।
\item \texttt{src} অ্যাট্রিবিউট না দেওয়া।
\item \texttt{src}-এর জায়গায় \texttt{href} লেখা।
\item ছবির ফাইলের নাম ভুল লেখা।
\item \texttt{.jpg} এর বদলে \texttt{.png} বা এক্সটেনশন ভুল দেওয়া।
\item ছবি অন্য ফোল্ডারে থাকলেও শুধু ফাইলের নাম লেখা।
\item \texttt{alt} অ্যাট্রিবিউট না দেওয়া।
\item \texttt{width} ও \texttt{height} উল্টো দেওয়া।
\item ছবিতে লিংক করতে \texttt{<img>} ট্যাগ \texttt{<a>} ট্যাগের বাইরে রাখা।
\end{itemize}
\end{multicols}


\noindent
\begin{minipage}[t]{.48\linewidth}\vspace{0pt}
\textbf{\color{chFourRose}ভুল কোড}\par
\begin{verbatim}
<image src="photo.jpg">
\end{verbatim}
\end{minipage}\hfill
\begin{minipage}[t]{.48\linewidth}\vspace{0pt}
\textbf{\color{chFourTeal}সঠিক কোড}\par
\begin{verbatim}
<img src="photo.jpg" alt="Photo">
\end{verbatim}
\end{minipage}

\subsection{পাঠের সারাংশ}



% Technical basis: MDN img. Correct modern syntax alongside legacy exam conventions.
\subsubsection{ছবির মাপ থেকে সঠিক কোড নির্বাচন}
প্রশ্নে উচ্চতা ও প্রস্থ যে ক্রমে লেখা থাকুক, \texttt{height} উচ্চতা এবং \texttt{width} প্রস্থ নির্দেশ করে। ছবির নাম বা পথ থাকে \texttt{src}-এ; \texttt{href} দিয়ে ছবি বসানো হয় না।

\par\noindent
\begin{minipage}[t]{.48\linewidth}\vspace{0pt}
\begin{chFourWholeExample}{HTML অ্যাট্রিবিউট: সংখ্যার সঙ্গে একক নয়}
\begin{verbatim}
<img src="badge.png"
     width="160" height="100"
     alt="Science club badge">
\end{verbatim}
\end{chFourWholeExample}
\end{minipage}\hfill
\begin{minipage}[t]{.48\linewidth}\vspace{0pt}
\begin{chFourWholeExample}{CSS-এ মাপ: px এককসহ}
\begin{verbatim}
<img src="badge.png"
     style="width:160px;height:100px"
     alt="Science club badge">
\end{verbatim}
\end{chFourWholeExample}
\end{minipage}\par
উভয় উদাহরণে প্রদর্শনের প্রস্থ ১৬০ ও উচ্চতা ১০০ CSS পিক্সেল। বর্তমান HTML-এ \texttt{width="160px"} সঠিক অ্যাট্রিবিউট মান নয়; \texttt{width="160"} লিখবে। পুরোনো প্রশ্নের বিকল্পে এমন এককযুক্ত মান থাকলে সেটিকে বর্তমান লেখার নিয়ম হিসেবে গ্রহণ করবে না; ট্যাগ, \texttt{src}, ফাইলের নাম এবং চাওয়া মাপ আলাদা করে বিচার করবে।


% Original applied explanation; see Chapter04/SOURCES.md.
\subsubsection{ছবির প্রদর্শনের মাপ ও ফাইলের মাপ আলাদা}
ক্যাম্পাসের বড় ছবির \texttt{width} কমালে পর্দায় ছোট দেখা যায়; তাতে মূল ফাইলের byte-সংখ্যা নিজে থেকে কমে না। ধীর সংযোগে দ্রুত দেখাতে উপযুক্ত মাপের ছবির ফাইল তৈরি করতে হয়। ছবির প্রস্থ ও উচ্চতার অনুপাত না মিলিয়ে দুটোই বদলালে ছবি বিকৃত হতে পারে।

ছবি দেখা না গেলে \texttt{src}-এর ফোল্ডার, নাম, বড়/ছোট হাতের অক্ষর এবং ফাইলের ধরন মিলিয়ে দেখবে। তথ্যবহুল ছবির \texttt{alt} তার প্রয়োজনীয় অর্থ দেয়; কেবল সাজসজ্জার ছবি হলে খালি \texttt{alt=""} ব্যবহার করা যায়।


\subsection{সংক্ষিপ্ত অনুশীলন}

% Each question keeps its answer beneath it in the same column.
\par\noindent
\begin{minipage}[t]{.48\linewidth}\vspace{0pt}
\chFourReviewQuestion{\textbf{HTML-এ ছবি যুক্ত করার ট্যাগ কোনটি?}}
\end{minipage}
\hfill
\begin{minipage}[t]{.48\linewidth}\vspace{0pt}
\chFourReviewQuestion{\textbf{\texttt{<img>} ট্যাগ কি কন্টেইনার ট্যাগ?}}
\end{minipage}
\par\smallskip

\par\noindent
\begin{minipage}[t]{.48\linewidth}\vspace{0pt}
\chFourReviewQuestion{\textbf{ছবির ফাইলের উৎস নির্দেশ করতে কোন অ্যাট্রিবিউট ব্যবহার হয়?}}
\end{minipage}
\hfill
\begin{minipage}[t]{.48\linewidth}\vspace{0pt}
\chFourReviewQuestion{\textbf{ছবির বিকল্প লেখা দিতে কোন অ্যাট্রিবিউট ব্যবহার হয়?}}
\end{minipage}
\par\smallskip

\par\noindent
\begin{minipage}[t]{.48\linewidth}\vspace{0pt}
\chFourReviewQuestion{\textbf{ছবির আকার নির্ধারণ করতে কোন অ্যাট্রিবিউট ব্যবহার হয়?}}
\end{minipage}
\hfill
\begin{minipage}[t]{.48\linewidth}\vspace{0pt}
\chFourReviewQuestion{\textbf{ছবিতে লিংক তৈরি করতে কী করতে হয়?}}
\end{minipage}
\par\smallskip

\par\noindent
\begin{minipage}[t]{.48\linewidth}\vspace{0pt}
\chFourReviewQuestion{\textbf{ওয়েবপেজে ছবির পরিচিত ফরম্যাট কী কী?}}
\end{minipage}
\hfill
\begin{minipage}[t]{.48\linewidth}\vspace{0pt}
\chFourReviewQuestion{\textbf{\texttt{src="images/logo.png"} মানে কী?}}
\end{minipage}
\par\smallskip

\par\noindent
\begin{minipage}[t]{.48\linewidth}\vspace{0pt}
\chFourReviewQuestion{\textbf{ধারকের মাপ অনুযায়ী ছবি নির্বাচনে কোন দুটি অ্যাট্রিবিউট বেশি গুরুত্বপূর্ণ?}}
\end{minipage}
\hfill
\begin{minipage}[t]{.48\linewidth}\vspace{0pt}
\chFourReviewQuestion{\textbf{প্রয়োজনের কাছাকাছি সময়ে ছবি লোড করার ইঙ্গিত দিতে কোন অ্যাট্রিবিউট ব্যবহার হয়?}}
\end{minipage}
\par\smallskip


\Needspace{10\baselineskip}
\subsection{পরীক্ষাভিত্তিক বহুনির্বাচনি প্রশ্ন}

\par\addvspace{4pt}\noindent
\begin{minipage}[t]{.48\linewidth}\vspace{0pt}
\begin{enumerate}[start=1,leftmargin=*,topsep=0pt,itemsep=0pt]
\item HTML-এ ছবি যুক্ত করার ট্যাগ কোনটি?
\par ক) \texttt{<ul>}
\par খ) \texttt{<ol>}
\par গ) \texttt{<img>}
\par ঘ) \texttt{<li>}
\end{enumerate}
\end{minipage}\hfill
\begin{minipage}[t]{.48\linewidth}\vspace{0pt}
\begin{enumerate}[start=2,leftmargin=*,topsep=0pt,itemsep=0pt]
\item \texttt{<img>} ট্যাগের ছবি উৎস অ্যাট্রিবিউট কোনটি?
\par ক) \texttt{href}
\par খ) \texttt{src}
\par গ) \texttt{alt}
\par ঘ) \texttt{link}
\end{enumerate}
\end{minipage}
\par
\par\addvspace{4pt}\noindent
\begin{minipage}[t]{.48\linewidth}\vspace{0pt}
\begin{enumerate}[start=3,leftmargin=*,topsep=0pt,itemsep=0pt]
\item ছবির আকার নিয়ন্ত্রণ করতে কোন অ্যাট্রিবিউট ব্যবহৃত হয়?
\par ক) \texttt{width} ও \texttt{height}
\par খ) \texttt{href} ও \texttt{target}
\par গ) \texttt{rowspan} ও \texttt{colspan}
\par ঘ) \texttt{type} ও \texttt{start}
\end{enumerate}
\end{minipage}\hfill
\begin{minipage}[t]{.48\linewidth}\vspace{0pt}
\begin{enumerate}[start=4,leftmargin=*,topsep=0pt,itemsep=0pt]
\item নিচের কোনটি সঠিক সিনট্যাক্স?
\par ক) \texttt{<image src="image.jpg">}
\par খ) \texttt{<img> src="image.jpg"}
\par গ) \texttt{<img src="image.jpg">}
\par ঘ) \texttt{<src img="image.jpg">}
\end{enumerate}
\end{minipage}
\par
\par\addvspace{4pt}\noindent
\begin{minipage}[t]{.48\linewidth}\vspace{0pt}
\begin{enumerate}[start=5,leftmargin=*,topsep=0pt,itemsep=0pt]
\item ছবি লোড না হলে যে লেখা দেখা যায়, সেটি কোন অ্যাট্রিবিউট দিয়ে দেওয়া হয়?
\par ক) \texttt{src}
\par খ) \texttt{alt}
\par গ) \texttt{width}
\par ঘ) \texttt{height}
\end{enumerate}
\end{minipage}\hfill
\begin{minipage}[t]{.48\linewidth}\vspace{0pt}
\begin{enumerate}[start=6,leftmargin=*,topsep=0pt,itemsep=0pt]
\item ছবিতে লিংক তৈরি করতে কোনটি সঠিক?
\par ক) \texttt{<a href="site.html"><img src="logo.png" alt="Logo"></a>}
\par খ) \texttt{<img href="site.html" src="logo.png">}
\par গ) \texttt{<a src="logo.png">site</a>}
\par ঘ) \texttt{<image link="site.html">}
\end{enumerate}
\end{minipage}
\par
\par\addvspace{4pt}\noindent
\begin{minipage}[t]{.48\linewidth}\vspace{0pt}
\begin{enumerate}[start=7,leftmargin=*,topsep=0pt,itemsep=0pt]
\item প্রয়োজনের কাছাকাছি সময়ে ছবি লোড করার ইঙ্গিত দিতে কোন অ্যাট্রিবিউট ব্যবহৃত হয়?
\par ক) \texttt{loading}
\par খ) \texttt{border}
\par গ) \texttt{align}
\par ঘ) \texttt{name}
\end{enumerate}
\end{minipage}\hfill
\begin{minipage}[t]{.48\linewidth}\vspace{0pt}
\begin{enumerate}[start=8,leftmargin=*,topsep=0pt,itemsep=0pt]
\item ধারকের মাপ অনুযায়ী ছবি নির্বাচনে কোন জোড়া অ্যাট্রিবিউট ব্যবহৃত হয়?
\par ক) \texttt{srcset} ও \texttt{sizes}
\par খ) \texttt{href} ও \texttt{target}
\par গ) \texttt{width} ও \texttt{height}
\par ঘ) \texttt{alt} ও \texttt{title}
\end{enumerate}
\end{minipage}
\par

\chFourMCQAnswer{১-গ, ২-খ, ৩-ক, ৪-গ, ৫-খ, ৬-ক, ৭-ক, ৮-ক}
